% ============================================================
%  Chapter 23: The Generative Spectral Database
%  Source: oxford/publications/prompt-based-spectral-database/
%          context-based-spectral-database.tex
% ============================================================

\epigraph{%
  The database is not a list of spectra.\\
  It is a traversal of phase space.%
}{}

\section{The Database as Phase Space}
\label{sec:ch23-phase-space}

A conventional spectral database stores pre-computed spectra as records,
indexed by molecular formula or identifier.
Query complexity is $O(N)$ where $N$ is the number of records.

In the partition framework, a molecule's mass spectrum is determined by its
partition cell at each fragmentation depth.
The ``database'' is the partition hierarchy of molecular space itself:
the S-entropy coordinate system $(\Sk, \St, \Se)$ is the address of any
molecular spectrum in the partition space.
There are no stored records; there is only the partition tree.

\begin{consequence}{O(k) Query Complexity}{Okquery}
  Query complexity of the generative spectral database is $O(k)$, where $k$
  is the partition depth of the target spectrum, independent of the total
  number of possible spectra $N$.
\end{consequence}

\begin{proof}[Proof sketch]
  Each query is a walk down the partition tree: at each depth level, one
  sub-cell is selected.
  A walk of depth $k$ requires $k$ steps.
  The total number of possible spectra $N = T(k_\mathrm{max}, d)$
  (composition inflation, Chapter~\ref{ch:inflation}) enters only in the
  maximum depth, not in the query time.
\end{proof}

\section{S-Entropy Coordinates as Database Address}
\label{sec:ch23-address}

The spectral address $(\Sk, \St, \Se)$ of a molecule uniquely identifies its
partition cell in the database:
\begin{itemize}
  \item $\Sk$ (knowledge entropy) indexes the chemical class and substructure.
  \item $\St$ (temporal entropy) indexes the chromatographic retention time.
  \item $\Se$ (partition entropy) indexes the mass analyser resolution depth.
\end{itemize}

Searching by $(\Sk, \St, \Se)$ is equivalent to walking the partition tree;
it is $O(k)$.

\section{Resolution Cascade and Bijection}
\label{sec:ch23-resolution}

A resolution cascade is the sequence of partition depths traversed from a
low-resolution ESI survey scan down to high-resolution MS$^n$ fragmentation.
Each step in the cascade:
\begin{enumerate}
  \item Narrows the $\Sk$ interval by the information content of the new
    measurement.
  \item Descends one level in the partition tree.
  \item Produces a new $(\Sk, \St, \Se)$ address for the candidate molecule.
\end{enumerate}

The ion--droplet bijection ensures that each ion droplet from ESI maps to a
unique path in the partition tree; there are no collisions at the resolution
depth of the instrument.

%%  SOURCE: integrate full derivations and panels from
%%          context-based-spectral-database.tex
%%  ADD: panel 1 (S-entropy space), panel 2 (analyser trajectories),
%%       panel 3 (cohesion clustering), panel 4 (bijection complexity).

\bigskip
\begin{tcolorbox}[colback=crownblue!5, colframe=crownblue!60!black,
                  title={\textbf{Chapter 23 Summary}}]
\begin{itemize}[noitemsep]
  \item The database is the partition hierarchy itself; spectra are
    partition-tree paths.
  \item Query complexity is $O(k)$ --- independent of database size $N$.
  \item S-entropy coordinates $(\Sk, \St, \Se)$ are the database address
    of any spectrum.
  \item The resolution cascade narrows the address at each MS$^n$ step.
\end{itemize}
\end{tcolorbox}
